\section{Appendix}

\subsection{Extractor Implementation Details}
\label{app:extractors}

\paragraph{Price extractor CSS selectors.}
The price extractor applies the following CSS selectors in order, returning on the
first match with $\conf = 0.9$:
\texttt{[data-testid='price']},
\texttt{.price},
\texttt{.listing-price},
\texttt{\#price},
\texttt{[class*='price']},
\texttt{[itemprop='price']},
\texttt{span[class*='Price']},
\texttt{div[class*='Price']}.
If no selector matches, the extractor falls back to a full-text regex with $\conf = 0.6$.
Extracted prices are normalized to lowercase with whitespace removed
(e.g., \texttt{\$2,400/mo}).

\paragraph{Availability patterns.}
The availability extractor matches the following regex patterns in order:
\begin{enumerate}
  \item \texttt{available\textbackslash s+(?:now|immediately)} $\to$ \texttt{available\_now}
  \item \texttt{available\textbackslash s+(\textbackslash w+\textbackslash s+\textbackslash d+,...)} $\to$ \texttt{available\_date}
  \item \texttt{(?:no longer|not)\textbackslash s+available} $\to$ \texttt{unavailable}
  \item \texttt{rented|leased|off\textbackslash s+market} $\to$ \texttt{off\_market}
  \item \texttt{move[- ]in\textbackslash s+(?:date\textbackslash s*:?\textbackslash s*)?...} $\to$ \texttt{available\_date}
\end{enumerate}
Structured selector matches yield $\conf = 0.9$; full-text matches yield $\conf = 0.65$.

\subsection{Candidate Collection: Craigslist}
\label{app:craigslist}

Craigslist listings are collected by scraping search results pages at
\texttt{https://sfbay.craigslist.org/search/apa?s=\{offset\}} with a 1.5-second
delay between requests. Individual listing URLs are extracted from
\texttt{li.cl-static-search-result a} elements whose \texttt{href} contains
\texttt{/d/} and ends in \texttt{.html}. This selector targets individual listing
permalinks and excludes search and category pages. Up to 100 unique URLs are collected
per run.

\subsection{Candidate Collection: Wikipedia}
\label{app:wikipedia}

Wikipedia candidates are collected via the MediaWiki API with the following
parameters:
\begin{verbatim}
action=query, list=recentchanges, rctype=edit,
rclimit=100, rcprop=title|timestamp|sizes,
rcnamespace=0, format=json
\end{verbatim}
The \texttt{User-Agent} header is set to
\texttt{FreshState-Research/1.0 (academic benchmark)}, as required by the
Wikimedia API terms of use. Omitting this header results in a 403 response.
URLs are constructed as \texttt{https://en.wikipedia.org/wiki/\{title\}}
with spaces replaced by underscores.

\subsection{Monitoring Infrastructure}
\label{app:infra}

The daily monitor is implemented in Python and scheduled via \texttt{cron}.
The cron expression \texttt{0 18 29,30,31 3 *} runs both domain monitors at
18:00 UTC (11:00 Pacific) on March 29, 30, and 31 2026. State is persisted
in JSON files (\texttt{monitor\_state\_apartment.json},
\texttt{monitor\_state\_product.json}) on the local filesystem between runs.
Logs are appended to \texttt{logs/monitor\_apartment.log} and
\texttt{logs/monitor\_product.log}.

Network errors and non-200 HTTP responses are logged and skipped; HTTP 410
responses are recorded as expiration events. A per-URL request timeout of
15 seconds is enforced to prevent stalled fetches from blocking the run.
